% R1_Cost_Definition.tex
% monogate Cost Theory — R1: SuperBEST Cost Definition and Basic Properties
% Defines expression trees, Cost_F, NaiveCost, and proves four foundational properties
% Date: 2026-04-20
% Author: Arturo R. Almaguer

\documentclass[12pt,a4paper]{article}
\usepackage{amsmath,amssymb,amsthm}
\usepackage{geometry}
\usepackage{booktabs}
\usepackage{array}
\usepackage{longtable}
\usepackage{xcolor}
\usepackage{hyperref}
\usepackage{listings}
\geometry{margin=2.5cm}

% ── Theorem environments ────────────────────────────────────────────────────
\newtheorem{theorem}{Theorem}[section]
\newtheorem{lemma}[theorem]{Lemma}
\newtheorem{corollary}[theorem]{Corollary}
\newtheorem{proposition}[theorem]{Proposition}
\theoremstyle{definition}
\newtheorem{definition}[theorem]{Definition}
\newtheorem{conjecture}[theorem]{Conjecture}
\newtheorem{remark}[theorem]{Remark}
\newtheorem{example}[theorem]{Example}
\newtheorem{observation}[theorem]{Observation}

% ── Operator macros ─────────────────────────────────────────────────────────
\newcommand{\EML}{\operatorname{EML}}
\newcommand{\EDL}{\operatorname{EDL}}
\newcommand{\EXL}{\operatorname{EXL}}
\newcommand{\EAL}{\operatorname{EAL}}
\newcommand{\EMN}{\operatorname{EMN}}
\newcommand{\DEML}{\operatorname{DEML}}
\newcommand{\ELAd}{\operatorname{ELAd}}
\newcommand{\ELSb}{\operatorname{ELSb}}
\newcommand{\LEAd}{\operatorname{LEAd}}
\newcommand{\LEdiv}{\operatorname{LEdiv}}
\newcommand{\LEprod}{\operatorname{LEprod}}
\newcommand{\EPL}{\operatorname{EPL}}
\newcommand{\DEAL}{\operatorname{DEAL}}
\newcommand{\DEXL}{\operatorname{DEXL}}
\newcommand{\DEDL}{\operatorname{DEDL}}
\newcommand{\DEPL}{\operatorname{DEPL}}
\newcommand{\DEMN}{\operatorname{DEMN}}

% ── Cost macros ──────────────────────────────────────────────────────────────
\newcommand{\Cost}{\operatorname{Cost}}
\newcommand{\NC}{\operatorname{NaiveCost}}
\newcommand{\SD}{\operatorname{SharingDiscount}}
\newcommand{\PB}{\operatorname{PatternBonus}}
\newcommand{\Fam}{\mathcal{F}}

% ── Lean listing style ───────────────────────────────────────────────────────
\lstdefinestyle{lean}{
  basicstyle=\ttfamily\small,
  keywordstyle=\bfseries,
  commentstyle=\itshape\color{gray},
  morekeywords={inductive,def,theorem,where,fun,match,with,if,then,else,%
                Type,Prop,Nat,iff,forall,exists,True,False,And,Or,Not,%
                structure,class,instance,open,namespace,end,import},
  literate={->}{{\(\to\)}}1 {=>}{{\(\Rightarrow\)}}1
           {>=}{{\(\geq\)}}1 {<=}{{\(\leq\)}}1
           {/\}{{\(\wedge\)}}1 {\/}{{\(\vee\)}}1,
  columns=flexible,
  keepspaces=true,
  frame=single,
  framerule=0.4pt,
  xleftmargin=8pt,
  xrightmargin=4pt,
  showstringspaces=false,
}

% ── Column types ─────────────────────────────────────────────────────────────
\newcolumntype{L}[1]{>{\raggedright\arraybackslash}p{#1}}
\newcolumntype{C}[1]{>{\centering\arraybackslash}p{#1}}

\title{SuperBEST Cost: Definition and Basic Properties\\[6pt]
       \large R1: Expression Trees, Cost Function, and Four Foundational Properties}
\author{Arturo R.\ Almaguer\\
\small monogate research \quad \texttt{almaguer1986@gmail.com}}
\date{April 2026}

% ════════════════════════════════════════════════════════════════════════════
\begin{document}
\maketitle

% ── Abstract ────────────────────────────────────────────────────────────────
\begin{abstract}
We give a precise mathematical definition of the \emph{SuperBEST cost} of an
arithmetic expression.  Section~\ref{sec:defs} introduces expression trees over
the 16-operator family $\Fam_{16}$, defines the minimal-DAG cost $\Cost_{\Fam}$
as a function on equivalence classes of arithmetic functions, and defines the
na\"{i}ve cost $\NC$ via the SuperBEST v3 unit-cost table.
Section~\ref{sec:props} proves four basic properties: non-negativity (P1),
the terminal characterisation of zero cost (P2), subadditivity under operator
application (P3), and algebraic invariance under pointwise equality of functions
(P4).  Section~\ref{sec:lean} sketches Lean~4 type signatures for the cost
function and all four properties.  These results form the foundational layer of
the General Theory of SuperBEST Cost developed in \texttt{Cost\_Theory.tex}.
\end{abstract}

\tableofcontents
\bigskip

% ════════════════════════════════════════════════════════════════════════════
\section{Definitions}
\label{sec:defs}
% ════════════════════════════════════════════════════════════════════════════

% ── 1.1 Expression Trees ────────────────────────────────────────────────────
\subsection{Expression Trees}

\begin{definition}[Terminal set]
\label{def:terminals}
Let $n, m \geq 0$.  The \emph{terminal set} is
\[
  \mathcal{T}
  \;=\;
  \{x_1,\ldots,x_n\}
  \;\cup\;
  \{c_1,\ldots,c_m\},
\]
where $x_1,\ldots,x_n$ are free \emph{variables} ranging over suitable subsets
of $\mathbb{R}$, and $c_1,\ldots,c_m$ are fixed real \emph{constants} (including
$0, 1, e, \pi$, etc.).
\end{definition}

\begin{definition}[SuperBEST operator family $\Fam_{16}$]
\label{def:operators}
The \emph{SuperBEST operator family} is the set
$\Fam_{16} = \{\EML, \DEML, \EXL, \EDL, \EAL, \EPL, \LEAd, \ELAd, \ELSb,
               \LEdiv, \LEprod, \DEAL, \DEXL, \DEDL, \DEPL, \DEMN\}$
of 16 binary operators, each mapping a pair of real inputs to a real output
on its natural domain.  The primary eleven operators are defined by:
\begin{align*}
  \EML(x,y)    &\;=\; e^{x} - \ln y, \\
  \DEML(x,y)   &\;=\; e^{-x} - \ln y, \\
  \EXL(x,y)    &\;=\; e^{x} \cdot \ln y, \\
  \EDL(x,y)    &\;=\; \frac{e^{x}}{\ln y}, \\
  \EAL(x,y)    &\;=\; e^{x} + \ln y, \\
  \EPL(x,y)    &\;=\; e^{x} + \ln y
                       \quad\text{(positive-domain variant of \EAL)}, \\
  \LEAd(x,y)   &\;=\; \ln(e^{x} + y)
                       \quad\text{(softplus-like)}, \\
  \ELAd(x,y)   &\;=\; e^{x} \cdot y
                       \quad\text{(since $e^{x+\ln y} = e^x \cdot y$)}, \\
  \ELSb(x,y)   &\;=\; e^{x} - y, \\
  \LEdiv(x,y)  &\;=\; \ln\!\left(\frac{e^{x}}{y}\right) = x - \ln y, \\
  \LEprod(x,y) &\;=\; \ln(e^{x} \cdot y) = x + \ln y.
\end{align*}
The remaining five operators $\DEAL, \DEXL, \DEDL, \DEPL, \DEMN$ are the
$e^{-x}$-variants of $\EAL, \EXL, \EDL, \EPL, \EML$ respectively, defined
analogously with $x$ replaced by $-x$.
\end{definition}

\begin{remark}[Functional overlap]
Several operators coincide on certain subsets of their domain.  In particular,
$\EAL(x,y) = \EPL(x,y)$ as abstract functions; the two names distinguish
contexts where the domain is restricted to positive arguments (\EPL) versus
the general real line (\EAL).  Similarly, $\LEdiv(x,y) = x - \ln y$, so
$\LEdiv$ implements subtraction when $y > 0$.  These overlaps are intentional:
they allow the SuperBEST search to select the operator whose unit cost is lowest
for the given sub-expression structure.
\end{remark}

\begin{definition}[Expression tree]
\label{def:expr-tree}
The set $\mathcal{E}(\mathcal{T}, \Fam)$ of \emph{expression trees} over
terminal set $\mathcal{T}$ and operator family $\Fam$ is defined inductively:
\begin{enumerate}
  \item \emph{(Terminal rule.)}
        Every $t \in \mathcal{T}$ is an expression tree, called a
        \emph{leaf} or \emph{terminal node}.
  \item \emph{(Operator rule.)}
        If $f \in \Fam$ and $A, B \in \mathcal{E}(\mathcal{T},\Fam)$ are
        expression trees, then $f(A, B)$ is an expression tree, called an
        \emph{internal node}.
\end{enumerate}
The \emph{size} $|E|$ of a tree $E$ is its total number of nodes (leaves and
internal nodes combined).  The \emph{depth} $d(E)$ is the length of the longest
root-to-leaf path.
\end{definition}

\begin{definition}[Semantic evaluation]
\label{def:eval}
Each expression tree $E \in \mathcal{E}(\mathcal{T},\Fam)$ defines a partial
function $\llbracket E \rrbracket : \mathbb{R}^n \to \mathbb{R}$ by structural
induction: a leaf $x_i$ evaluates to the $i$-th coordinate; a leaf $c_j$
evaluates to the constant $c_j \in \mathbb{R}$; an internal node $f(A,B)$
evaluates to $f(\llbracket A\rrbracket(\mathbf{x}),
                 \llbracket B\rrbracket(\mathbf{x}))$
wherever both sub-expressions and $f$ are defined.  Two trees $E_1, E_2$ are
\emph{semantically equivalent}, written $E_1 \equiv E_2$, when
$\llbracket E_1 \rrbracket = \llbracket E_2 \rrbracket$ pointwise on
a common domain of full measure.
\end{definition}

% ── 1.2 DAG Semantics ───────────────────────────────────────────────────────
\subsection{DAG Semantics}

\begin{definition}[Expression DAG over $\Fam$]
\label{def:dag}
An \emph{expression DAG} (directed acyclic graph) over $\Fam$ is a finite DAG
$D = (V, A)$ together with:
\begin{itemize}
  \item a labelling $\ell : V \to \mathcal{T} \cup \Fam$ such that every
        \emph{source} node (in-degree 0) is labelled by an element of
        $\mathcal{T}$, and every internal node $v$ (in-degree 2) is labelled
        by some $f_v \in \Fam$;
  \item a designated \emph{root} $r \in V$ (out-degree 0) giving the output.
\end{itemize}
Shared sub-expressions appear as nodes with out-degree $\geq 2$.  The DAG
computes a function $\llbracket D \rrbracket : \mathbb{R}^n \to \mathbb{R}$
defined by evaluating each node in topological order.
The \emph{node count} of $D$ is $|V_{\mathrm{int}}|$, the number of internal
(non-leaf) nodes in $D$.
\end{definition}

% ── 1.3 Cost Function ───────────────────────────────────────────────────────
\subsection{The Cost Function}

\begin{definition}[SuperBEST cost]
\label{def:cost}
Let $E \in \mathcal{E}(\mathcal{T},\Fam_{16})$ be an expression tree computing
the function $\llbracket E \rrbracket$.  The \emph{SuperBEST cost} (or simply
\emph{cost}) of $E$ is
\[
  \Cost_{\Fam}(E)
  \;=\;
  \min\bigl\{
    |V_{\mathrm{int}}(D)|
    \;\bigm|\;
    D \text{ is a DAG over } \Fam_{16},\;
    \llbracket D \rrbracket \equiv \llbracket E \rrbracket
  \bigr\},
\]
where the minimum ranges over all finite expression DAGs over $\Fam_{16}$ that
compute the same function as $E$ pointwise on a domain of full measure.  We
write $\Cost(E)$ when $\Fam = \Fam_{16}$ is clear from context.
\end{definition}

\begin{remark}[Cost is a property of functions, not trees]
\label{rem:cost-class}
Because the minimum in Definition~\ref{def:cost} depends only on the function
$\llbracket E \rrbracket$ and not on the syntactic structure of $E$, we have
$\Cost(E_1) = \Cost(E_2)$ whenever $E_1 \equiv E_2$.  The cost function
therefore factors through the equivalence classes of $\equiv$: it is a
well-defined function on (partial) functions $\mathbb{R}^n \to \mathbb{R}$
that are computable by some DAG over $\Fam_{16}$.
\end{remark}

\begin{remark}[Existence of the minimum]
The minimum exists because $|V_{\mathrm{int}}(D)| \geq 0$ for every DAG, the
set of achievable node counts is a non-empty subset of $\mathbb{N}_0$ (the
original tree $E$ itself gives a finite DAG), and every non-empty subset of
$\mathbb{N}_0$ has a minimum.
\end{remark}

% ── 1.4 NaiveCost ───────────────────────────────────────────────────────────
\subsection{The Na\"{i}ve Cost Function}

\begin{definition}[SuperBEST v3 unit costs]
\label{def:unit-costs}
The SuperBEST v3 routing table assigns a positive integer \emph{unit cost}
$c_{\mathrm{op}} \in \mathbb{Z}_{>0}$ to each primitive arithmetic operation:
\begin{center}
\begin{tabular}{@{}llc@{}}
\toprule
\textbf{Operation} & \textbf{Notation} & \textbf{Unit cost $c_{\mathrm{op}}$} \\
\midrule
Exponential $e^x$            & \texttt{exp}   & 1 \\
Natural logarithm $\ln x$    & \texttt{ln}    & 1 \\
Division $x / y$             & \texttt{div}   & 1 \\
Negation $-x$                & \texttt{neg}   & 2 \\
Reciprocal $1/x$             & \texttt{recip} & 2 \\
Multiplication $x \cdot y$   & \texttt{mul}   & 2 \\
Subtraction $x - y$          & \texttt{sub}   & 2 \\
Exponentiation $x^n$         & \texttt{pow}   & 3 \\
Addition $x+y$ (pos.\ domain)& \texttt{add}$^+$ & 3 \\
Addition $x+y$ (general)     & \texttt{add}   & 11 \\
\bottomrule
\end{tabular}
\end{center}
These counts are optimal in $\Fam_{16}$: no DAG over $\Fam_{16}$ computing
the given primitive with fewer internal nodes exists.
\end{definition}

\begin{definition}[Na\"{i}ve Cost]
\label{def:NC}
For any arithmetic expression $E$, the \emph{na\"{i}ve cost} is
\[
  \NC(E)
  \;=\;
    1 \cdot n_{\exp}
  + 1 \cdot n_{\ln}
  + 1 \cdot n_{\mathrm{div}}
  + 2 \cdot n_{\mathrm{neg}}
  + 2 \cdot n_{\mathrm{rec}}
  + 2 \cdot n_{\mathrm{mul}}
  + 2 \cdot n_{\mathrm{sub}}
  + 3 \cdot n_{\mathrm{pow}}
  + 3 \cdot n_{\mathrm{add}},
\]
where $n_{\mathrm{op}}$ counts the number of occurrences of each primitive
operation in the expression tree $E$ (without sharing).  Leaf nodes (variables
and constants) contribute~$0$.
\end{definition}

\begin{remark}[Relationship to $\Cost$]
The na\"{i}ve tree $T_{\mathrm{naive}}(E)$, obtained by substituting each
primitive in $E$ with its SuperBEST v3 optimal sub-tree independently, achieves
exactly $\NC(E)$ internal nodes.  Hence $\Cost(E) \leq \NC(E)$ for all $E$
(Theorem T34 of \texttt{Cost\_Theory.tex}).  The gap $\NC(E) - \Cost(E)$ is
exactly the sum of all sharing discounts and pattern bonuses.
\end{remark}

% ════════════════════════════════════════════════════════════════════════════
\section{Four Basic Properties}
\label{sec:props}
% ════════════════════════════════════════════════════════════════════════════

We now establish four fundamental properties of $\Cost$.  Throughout this
section, $E$, $A$, $B$ denote expression trees in
$\mathcal{E}(\mathcal{T}, \Fam_{16})$, and $\mathrm{op}$ denotes an arbitrary
element of $\Fam_{16}$.

% ── P1: Non-negativity ──────────────────────────────────────────────────────
\subsection{Property P1: Non-negativity}

\begin{proposition}[P1 --- Non-negativity]
\label{prop:P1}
For every expression tree $E \in \mathcal{E}(\mathcal{T},\Fam_{16})$,
\[
  \Cost(E) \;\geq\; 0.
\]
\end{proposition}

\begin{proof}
By Definition~\ref{def:cost}, $\Cost(E)$ is the minimum of the set
$S = \{|V_{\mathrm{int}}(D)| \mid D \text{ is a valid DAG for } E\}$.
Every element of $S$ is of the form $|V_{\mathrm{int}}(D)|$, a cardinality of
a finite set, which satisfies $|V_{\mathrm{int}}(D)| \geq 0$.  Therefore every
element of $S$ is non-negative, and $\Cost(E) = \min S \geq 0$.
\end{proof}

% ── P2: Zero cost iff terminal ───────────────────────────────────────────────
\subsection{Property P2: Zero Cost if and only if Terminal}

\begin{proposition}[P2 --- Terminal Characterisation]
\label{prop:P2}
For every $E \in \mathcal{E}(\mathcal{T},\Fam_{16})$:
\[
  \Cost(E) = 0
  \quad\iff\quad
  E \text{ is a terminal (leaf) node}.
\]
\end{proposition}

\begin{proof}
\textbf{($\Leftarrow$)}
If $E = t$ for some $t \in \mathcal{T}$, then the single-node DAG consisting
only of the leaf $t$ computes $\llbracket E \rrbracket$ with zero internal
nodes.  Hence $\Cost(E) \leq 0$, and by P1, $\Cost(E) = 0$.

\textbf{($\Rightarrow$)}
Suppose $E$ is not a terminal.  Then $E = f(A, B)$ for some $f \in \Fam_{16}$
and sub-trees $A$, $B$.  Any DAG $D$ computing $\llbracket E \rrbracket
= f(\llbracket A \rrbracket, \llbracket B \rrbracket)$ must include at least
one internal node: a node labelled $f$ (or some operator that produces the same
function as $f$ on the given sub-expression values).  Therefore
$|V_{\mathrm{int}}(D)| \geq 1$ for every valid DAG $D$, so $\Cost(E) \geq 1
> 0$.
\end{proof}

\begin{remark}
The forward implication is the non-trivial direction: it uses the fact that no
operator in $\Fam_{16}$ can be ``compiled away'' to zero nodes.  Every $f \in
\Fam_{16}$ computes a function strictly different from the identity on at least
one sub-expression shape, so at least one physical node is required.
\end{remark}

% ── P3: Subadditivity ───────────────────────────────────────────────────────
\subsection{Property P3: Subadditivity Under Operator Application}

\begin{proposition}[P3 --- Subadditivity]
\label{prop:P3}
For all $\mathrm{op} \in \Fam_{16}$ and all
$A, B \in \mathcal{E}(\mathcal{T},\Fam_{16})$:
\[
  \Cost(\mathrm{op}(A, B))
  \;\leq\;
  \Cost(A) + \Cost(B) + 1.
\]
\end{proposition}

\begin{proof}
Let $D_A$ be an optimal DAG for $A$ (so $|V_{\mathrm{int}}(D_A)| = \Cost(A)$)
and let $D_B$ be an optimal DAG for $B$ (so $|V_{\mathrm{int}}(D_B)| = \Cost(B)$).
Construct the DAG $D$ as follows:
\begin{enumerate}
  \item Take the disjoint union of $D_A$ and $D_B$, identifying any shared leaf
        nodes (variables or constants that appear in both).  This step does not
        increase the internal node count.
  \item Add a single new internal node $v^*$ labelled $\mathrm{op}$, with
        incoming edges from the roots of $D_A$ and $D_B$ respectively, and
        declare $v^*$ to be the root of $D$.
\end{enumerate}
The resulting graph $D$ is a valid DAG (no cycles are introduced, and
$v^*$ is new).  By construction,
\[
  \llbracket D \rrbracket
  = \mathrm{op}(\llbracket D_A \rrbracket, \llbracket D_B \rrbracket)
  = \mathrm{op}(\llbracket A \rrbracket, \llbracket B \rrbracket)
  = \llbracket \mathrm{op}(A,B) \rrbracket,
\]
so $D$ is a valid DAG for $\mathrm{op}(A,B)$.  Its internal node count is
\[
  |V_{\mathrm{int}}(D)|
  = |V_{\mathrm{int}}(D_A)| + |V_{\mathrm{int}}(D_B)| + 1
  = \Cost(A) + \Cost(B) + 1.
\]
Since $\Cost(\mathrm{op}(A,B))$ is the \emph{minimum} over all valid DAGs:
\[
  \Cost(\mathrm{op}(A,B))
  \;\leq\;
  |V_{\mathrm{int}}(D)|
  = \Cost(A) + \Cost(B) + 1. \qedhere
\]
\end{proof}

\begin{remark}[When the bound is tight]
The bound $\Cost(A) + \Cost(B) + 1$ is achieved whenever no algebraic
identity allows $\mathrm{op}(A,B)$ to be compressed further.  For example,
$\Cost(\EML(A, B)) = \Cost(A) + \Cost(B) + 1$ when $A$ and $B$ are
independent variables.  The inequality can be \emph{strict} when a pattern
bonus applies: e.g.\ $\LEprod(x, e^y) = x + y$ compresses the na\"{i}ve
$1 + 1 + 1 = 3$ to $\Cost = 0$ (pure addition of terminals).
\end{remark}

\begin{corollary}[Upper bound by structural induction]
\label{cor:UB-inductive}
For any expression tree $E$ with $k$ internal nodes, $\Cost(E) \leq k$.
\end{corollary}

\begin{proof}
By induction on $k$: the base case $k = 0$ is P2 ($\Cost = 0$).  For $k \geq 1$,
write $E = \mathrm{op}(A,B)$.  By the inductive hypothesis,
$\Cost(A) \leq |A|_{\mathrm{int}}$ and $\Cost(B) \leq |B|_{\mathrm{int}}$.
By P3, $\Cost(E) \leq |A|_{\mathrm{int}} + |B|_{\mathrm{int}} + 1 = k$.
\end{proof}

% ── P4: Algebraic Invariance ─────────────────────────────────────────────────
\subsection{Property P4: Algebraic Invariance}

\begin{proposition}[P4 --- Algebraic Invariance]
\label{prop:P4}
If $E_1, E_2 \in \mathcal{E}(\mathcal{T},\Fam_{16})$ satisfy
$E_1 \equiv E_2$ (i.e.\ $\llbracket E_1 \rrbracket = \llbracket E_2 \rrbracket$
pointwise on a domain of full measure), then
\[
  \Cost(E_1) = \Cost(E_2).
\]
\end{proposition}

\begin{proof}
By Definition~\ref{def:cost}, $\Cost(E_i)$ depends only on the function
$\llbracket E_i \rrbracket$: it is the minimum number of internal nodes over
all DAGs $D$ satisfying $\llbracket D \rrbracket \equiv \llbracket E_i
\rrbracket$.  Since $\llbracket E_1 \rrbracket = \llbracket E_2 \rrbracket$,
the two minimisation problems are identical:
\begin{align*}
  \Cost(E_1)
  &= \min\bigl\{|V_{\mathrm{int}}(D)| \;\bigm|\;
      \llbracket D \rrbracket \equiv \llbracket E_1 \rrbracket\bigr\} \\
  &= \min\bigl\{|V_{\mathrm{int}}(D)| \;\bigm|\;
      \llbracket D \rrbracket \equiv \llbracket E_2 \rrbracket\bigr\} \\
  &= \Cost(E_2). \qedhere
\end{align*}
\end{proof}

\begin{remark}[Practical consequence]
P4 is the formal justification for all algebraic simplifications used in the
SuperBEST search.  Replacing $e^{x} \cdot e^{y}$ by $e^{x+y}$, or
$\ln(x/y)$ by $\ln x - \ln y$, does not change the cost because both
expressions denote the same function.  Cost optimisation may therefore proceed
in any algebraically equivalent representation.
\end{remark}

\begin{remark}[P4 and the cost of constants]
A special case: if $E$ is an expression all of whose sub-expressions evaluate
to constants (e.g.\ $\ln(e^2)$), then $E \equiv c$ for some $c \in \mathbb{R}$,
and by P4, $\Cost(E) = \Cost(c) = 0$.  This is constant folding, one of the
two mechanisms behind the Sharing Discount in T38.
\end{remark}

% ── Summary table ────────────────────────────────────────────────────────────
\subsection{Summary of Basic Properties}

\begin{table}[h]
\centering
\caption{The four basic properties of SuperBEST cost.}
\label{tab:props}
\smallskip
\begin{tabular}{@{}clL{7.5cm}@{}}
\toprule
\textbf{Label} & \textbf{Name} & \textbf{Statement} \\
\midrule
P1 & Non-negativity
   & $\Cost(E) \geq 0$ for all $E$. \\[4pt]
P2 & Terminal characterisation
   & $\Cost(E) = 0 \iff E$ is a leaf node. \\[4pt]
P3 & Subadditivity
   & $\Cost(\mathrm{op}(A,B)) \leq \Cost(A) + \Cost(B) + 1$. \\[4pt]
P4 & Algebraic invariance
   & $E_1 \equiv E_2 \Rightarrow \Cost(E_1) = \Cost(E_2)$. \\
\bottomrule
\end{tabular}
\end{table}

% ════════════════════════════════════════════════════════════════════════════
\section{Lean 4 Type Signatures}
\label{sec:lean}
% ════════════════════════════════════════════════════════════════════════════

This section sketches Lean~4 type signatures for the definitions and properties
of Sections~\ref{sec:defs}--\ref{sec:props}.  These are \emph{type-level stubs}
only: the proofs are omitted.  A full Lean formalisation is future work.

\begin{lstlisting}[style=lean, caption={Expression tree and operator types.}]
-- The family of SuperBEST operators
inductive Operator : Type where
  | EML  | DEML | EXL  | EDL  | EAL  | EPL
  | LEAd | ELAd | ELSb | LEdiv | LEprod
  | DEAL | DEXL | DEDL | DEPL | DEMN
  deriving DecidableEq, Repr

-- Expression tree parameterised by a terminal type alpha
inductive Expr (alpha : Type) : Type where
  | terminal : alpha -> Expr alpha
  | node     : Operator -> Expr alpha -> Expr alpha -> Expr alpha
  deriving Repr

-- Predicate: is a leaf node?
def Expr.isTerminal {alpha : Type} : Expr alpha -> Prop
  | terminal _ => True
  | node _ _ _ => False
\end{lstlisting}

\begin{lstlisting}[style=lean, caption={Cost function and NaiveCost.}]
-- Unit cost of each operator (SuperBEST v3)
def unitCost : Operator -> Nat
  | .EML | .DEML | .EXL  | .EDL  | .EAL
  | .EPL | .DEAL | .DEXL | .DEDL | .DEPL | .DEMN => 1
  -- (actual values from Table 1; simplified here for illustration)
  | _ => 1

-- NaiveCost: sum of unit costs over all internal nodes
def naiveCost {alpha : Type} : Expr alpha -> Nat
  | Expr.terminal _   => 0
  | Expr.node op a b  => unitCost op + naiveCost a + naiveCost b

-- SuperBEST cost (axiomatic; minimisation over all equivalent DAGs)
-- In a full formalisation this would be defined via a DAG type.
noncomputable def cost : Expr Real -> Nat :=
  fun e => Nat.find (cost_exists e)   -- where cost_exists : exists n, ...
\end{lstlisting}

\begin{lstlisting}[style=lean, caption={Type signatures for the four basic properties.}]
-- P1: Non-negativity (trivially true for Nat, stated for documentation)
theorem cost_nonneg (e : Expr Real) : cost e >= 0 :=
  Nat.zero_le _

-- P2: Zero cost iff terminal
theorem cost_zero_iff_terminal (e : Expr Real) :
    cost e = 0 <-> e.isTerminal := by
  sorry  -- proof omitted; see Proposition P2

-- P3: Subadditivity under operator application
theorem cost_subadditive
    (op : Operator) (a b : Expr Real) :
    cost (Expr.node op a b) <= cost a + cost b + 1 := by
  sorry  -- proof omitted; see Proposition P3

-- P4: Algebraic invariance
-- Requires a semantic evaluation function eval : Expr Real -> (Real -> Real)
theorem cost_algebraic_invariant
    (e1 e2 : Expr Real)
    (h : forall x, eval e1 x = eval e2 x) :
    cost e1 = cost e2 := by
  sorry  -- proof omitted; see Proposition P4
\end{lstlisting}

\begin{remark}[Status of the Lean formalisation]
P1 is vacuous in Lean because \texttt{cost} returns a \texttt{Nat}, which is
always non-negative by type.  P2 requires a formal definition of the DAG
minimisation; the stub uses \texttt{Nat.find} to assert existence.  P3 and P4
follow from the minimisation semantics and are marked \texttt{sorry} pending a
full DAG formalisation.  The four signatures are provided as a bridge to a
future mechanised proof in Lean~4 or Mathlib.
\end{remark}

% ════════════════════════════════════════════════════════════════════════════
\section*{Acknowledgments}
% ════════════════════════════════════════════════════════════════════════════

This document is part of the monogate cost theory sprint (2026-04-20) and
constitutes the foundational layer (R1) underlying the results in
\texttt{Cost\_Theory.tex} (T34--T39).  Definitions and notation follow
\texttt{Cost\_Theory.tex} exactly.

\bigskip
\noindent\textit{Almaguer, A.R.\ (2026).
SuperBEST Cost: Definition and Basic Properties (R1).
monogate research. \url{https://monogate.org}}

\end{document}
